% Generated draft: requires mathematical and source review.
\documentclass[11pt]{article}
\usepackage[margin=25mm]{geometry}
\usepackage{fontspec}
\setmainfont{Latin Modern Roman}
\usepackage{amsmath,amssymb,mathtools,unicode-math}
\setmathfont{Latin Modern Math}
\usepackage{enumitem,needspace,xurl,hyperref}
\setlength{\parindent}{0pt}
\setlength{\parskip}{5pt}
\setlength{\emergencystretch}{3em}
\newcommand{\N}{\mathbb{N}}
\newcommand{\Z}{\mathbb{Z}}
\newcommand{\Q}{\mathbb{Q}}
\newcommand{\R}{\mathbb{R}}
\newcommand{\C}{\mathbb{C}}
\newcommand{\F}{\mathbb{F}}
\newcommand{\A}{\mathbb{A}}
\newcommand{\PP}{\mathbb{P}}
\newcommand{\E}{\mathbb{E}}
\newcommand{\eps}{\varepsilon}
\newcommand{\dd}{\,\mathrm{d}}
\begin{document}
\section*{UnsolvedMath ID: 6600001}
\textbf{Source problem code:} catalogue entry
\subsection{Definitions and mathematical statement}
Let \(D=\N_{\ge0}\) for a one-sided sequence or \(D=\Z\) for a two-sided sequence, and let \(s:D\to\mathcal A\). For a finite offset pattern \(\tau=(0=\tau_0<\tau_1<\cdots<\tau_{k-1})\) of integers, define
\[
\mathcal P_s(\tau)=\{(s(n+\tau_0),\ldots,s(n+\tau_{k-1})):n\in D\},
\qquad
p_s^*(k)=\sup_\tau|\mathcal P_s(\tau)|.
\]
The offsets are nonnegative, so every sampled index is in \(D\). A sequence is pattern Sturmian in the usual maximal-pattern-complexity sense when
\[
\forall k\ge1,\qquad p_s^*(k)=2k.
\]
In particular exactly two symbols occur, by \(k=1\). This uses arbitrary finite offset sets, not just consecutive blocks as in ordinary Sturmian complexity.

The problem is to classify all such \(s\) separately for \(D=\N_{\ge0}\) and \(D=\Z\), including whichever recurrent and nonrecurrent sequences occur in either domain. A classification must specify all parameters and exceptional cases of the proposed constructions, and how shift or symbol relabeling is treated if equivalence classes are used.

The supplied research summary asserts a complete resolution as of 2025: recurrent sequences are described by irrational-circle two-interval codings and nearly simple Toeplitz subshifts, while nonrecurrent sequences are described as very close to constant or by nonrecurrent rotation codings. It supplies no precise reference, construction definitions, or parameter restrictions. Recurrent is not synonymous with two-sided, and nonrecurrent is not synonymous with one-sided; the summary's parenthetical identifications do not replace the two classification tasks. Very close to constant is not a mathematical class without a specified sparsity rule. The claimed complete classification is therefore retained as unverified metadata rather than an explicit exhaustive theorem.
\subsection{Short English statement}
Classify one-sided and two-sided sequences whose maximal complexity over every \(k\)-position pattern is exactly \(2k\). The supplied solved assessment lacks a precise classification reference and conflates index direction with recurrence.
\subsection{Sources}
\begin{itemize}
\item Local UnsolvedMath registry, permanent ID \(6600001\), exact source code \texttt{catalogue entry}, ``D. Damanik: Quantum Mechanics and Quasicrystals --- Problem.'' Proposed year 2016.
\item \emph{Open Problems and Conjectures Related to the Theory of Mathematical Quasicrystals} (2016), source-order Problem 1, D. Damanik section ``Quantum Mechanics and Quasicrystals''; \url{https://arxiv.org/abs/1604.06280}. Source-TeX extraction; recorded access 2026-07-29.
\item Supplied research summary claims a complete classification as of 2025 but supplies no authors, title, publication, or precise theorem for that claim.
\end{itemize}
\end{document}
